\subsection{Generalized influence implies hereditary full AA}
\label{ian:generalized-influence-bridge}

Status: proved.

Throughout, \(D_i f=(f(x_{-i},1)-f(x_{-i},-1))/2\),
\(I_i^\mu(f)=\E_\mu(D_i f)^2\),
\[
 \mathcal I_\mu(f)=\sum_i I_i^\mu(f),
 \qquad
 G_\mu(f,S)=\Var_\mu\bigl(\E_\mu[f\mid X_S]\bigr),
\]
the latter being the reveal gain of Section~\ref{tht:reveal}.
All polynomial bounds on functions are bounds on the entire cube.

\begin{definition}[Generalized AA]
\label{ian:gamma-definition}
Let \(\Lambda\) be a class of laws on finite cubes. We say that
\(\Gamma^{\mathrm{rand}}(\Lambda)\) holds if there are constants
\[
 c_L,C_L,c_G,C_G,B_G,K_G>0
\]
such that the following assertion holds for every \(n,d\ge1\), every
\(\mu\in\Lambda\) on \(\{-1,1\}^n\), every \(v\in(0,1]\), and every
multilinear \(f:\{-1,1\}^n\to[-1,1]\) of degree at most \(d\) with
\(\Var_\mu f\ge v\). At least one of the following holds:
\[
 \max_i I_i^\mu(f)\ge c_L(v/d)^{C_L};
\]
or, with
\[
 \ell=\min\{n,\lceil B_G(d/v)^{K_G}\rceil\},
\]
a uniformly random \(\ell\)-subset \(S\), independent of \(X\sim\mu\),
satisfies
\[
 \E_S G_\mu(f,S)\ge c_G(v/d)^{C_G}.
\]
The assertion \(\Gamma^{\mathrm{set}}(\Lambda)\) replaces the second
alternative by the existence of a fixed \(S\subseteq[n]\), with
\(|S|\le\ell\), satisfying the same lower bound without \(\E_S\).
All six constants are chosen before \(\mu,n,d,v,f\).
\end{definition}

Some \(\ell\)-set attains at least the average, so
\(\Gamma^{\mathrm{rand}}(\Lambda)\) implies
\(\Gamma^{\mathrm{set}}(\Lambda)\) with the same constants.

\begin{lemma}[Potentials under arbitrary input laws]
\label{ian:query-potentials}
For every law \(\mu\) on \(\{-1,1\}^n\) and every multilinear
\(f:\{-1,1\}^n\to[-1,1]\) of degree at most \(d\ge1\),
\[
 \sum_i(D_i f(x))^2\le2d^2
 \quad\text{for every }x,
 \qquad
 \mathcal I_\mu(f)\le2d^2.
\]
For every \(A\subseteq[n]\) and every \(j\notin A\),
\[
 \E_{X_A} I_j^{\mu\mid X_A}(f\mid X_A)=I_j^\mu(f).
\]
Here and below expectations over pinnings omit zero-probability
assignments and coordinates are relabelled when necessary. Consequently,
\[
 \E_{X_A}\mathcal I_{\mu\mid X_A}(f\mid X_A)
 =\mathcal I_\mu(f)-\sum_{i\in A}I_i^\mu(f),
 \tag{1}\label{ian:influence-decrement}
\]
and
\[
 \E_{X_A}\Var_{\mu\mid X_A}(f\mid X_A)
 =\Var_\mu f-G_\mu(f,A).
 \tag{2}\label{ian:variance-decrement}
\]
Thus both potentials are supermartingales along an adaptive sequence
of coordinate queries.
\end{lemma}

\begin{proof}
Fix \(x\), write \(a_i=x_iD_i f(x)\), and fix \(A\subseteq[n]\).
Independently flip each coordinate in \(A\) with probability \(t\),
leaving the others equal to \(x\). Then \(h_A(t)=\E f(X(t))\) is a
polynomial of degree at most \(d\), bounded by one on \([0,1]\), with
\[
 h_A'(0)=\sum_{i\in A}\bigl(f(x^i)-f(x)\bigr)
        =-2\sum_{i\in A}a_i.
\]
Markov brothers' inequality on \([0,1]\),
\(\|p'\|\le2\deg(p)^2\|p\|\), gives
\(\lvert\sum_{i\in A}a_i\rvert\le d^2\).
Taking \(A\) to be the set of positive, then of negative, \(a_i\) and
using \(|D_i f|\le1\),
\[
 \sum_i(D_i f(x))^2
 \le\sum_i|a_i|\le2d^2.
\]
Integration proves the total-influence bound. Girish Lemma~2.11, p.~9,
and the proof of (5) in
Section~\ref{influence-recovery:exact-potential-identity-and-dimension-free-budget}
give the constant \(4d^2\) in these conventions.

For \(j\notin A\), \(D_j(f\mid X_A)=(D_j f)\mid X_A\), and the tower
property applied to its square proves the influence identity, with no
independence assumption. Summing over \(j\notin A\) proves
\eqref{ian:influence-decrement}; total variance proves
\eqref{ian:variance-decrement}. Both apply conditionally at every
positive-probability transcript of an adaptive tree.
\end{proof}

\begin{theorem}[Theorem A: the bridge]
\label{ian:bridge}
Suppose that \(\Lambda\) is pinning-closed and that either
\(\Gamma^{\mathrm{rand}}(\Lambda)\) or
\(\Gamma^{\mathrm{set}}(\Lambda)\) holds. Then \(\Lambda\) has hereditary
full AA: full AA in the sense of
\path{synthesis/threshold_theory/01_target_and_conclusions.tex},
Section~\ref{tht:target}, uniformly over all members of \(\Lambda\) and hence over all their positive-probability pinnings, as described after Proposition~\ref{tht:reveal-prop}. (For \(\varepsilon\ge1\) the constant output \(1/2\) suffices.)

More generally, let \(\mu\in\Lambda\), let
\(Q:\{-1,1\}^n\to[0,1]\) be a multilinear polynomial of degree at most
\(d\ge1\), and let \(0<\varepsilon<1\). Put
\[
 a=\varepsilon^2/4,\qquad
 \alpha=c_L(a/d)^{C_L},\qquad
 \gamma=c_G(a/d)^{C_G},\qquad
 L=\lceil B_G(d/a)^{K_G}\rceil .
\]
There is a deterministic tree with outputs in \([0,1]\), mean absolute
error at most \(\varepsilon\) under \(\mu\), and depth at most
\[
 \left\lceil
 \frac{2}{\varepsilon}
 \left(\frac{2d^2}{\alpha}+\frac{L}{\gamma}\right)
 \right\rceil.
 \tag{3}\label{ian:bridge-depth}
\]
\end{theorem}

\begin{proof}
By the observation after Definition~\ref{ian:gamma-definition}, we may
assume \(\Gamma^{\mathrm{set}}(\Lambda)\).
First build a tree without a depth cutoff. At a positive-probability
transcript, let \(\nu\) be the residual law, which lies in \(\Lambda\)
by pinning closure, and \(q\) the restricted polynomial. If no
coordinates remain, or if \(\Var_\nu q<a\), stop and output
\(\E_\nu q\). Otherwise apply \(\Gamma^{\mathrm{set}}\) to \(q\) with
variance parameter \(a\). If some remaining coordinate has
\(I_i^\nu(q)\ge\alpha\), query one such coordinate (a Local step).
Otherwise query a set \(S\) of at most \(L\) remaining coordinates with
\(G_\nu(q,S)\ge\gamma\) (a Global step); \(S\neq\varnothing\) because
its gain is positive. Every step queries a new coordinate, so the tree
has depth at most \(n\).

Let \(N_L,N_G\) count the Local and Global steps, and let \(K\) be the
total number of queries. By \eqref{ian:influence-decrement}, each Local
step decreases expected total influence by at least \(\alpha\) and no
step increases it; by \eqref{ian:variance-decrement}, each Global step
decreases expected conditional variance by at least \(\gamma\) and no
step increases it. Telescoping these finite sums and using
Lemma~\ref{ian:query-potentials},
\[
 \alpha\,\E N_L\le\mathcal I_\mu(Q)\le2d^2,
 \qquad
 \gamma\,\E N_G\le\Var_\mu Q\le1.
\]
Since \(K\le N_L+LN_G\),
\[
 \E K\le H:=\frac{2d^2}{\alpha}+\frac{L}{\gamma}.
 \tag{4}\label{ian:expected-query-bound}
\]

Write \(D_*\) for the output of this tree. Every leaf has conditional
variance at most \(a\), so conditional Cauchy--Schwarz gives
\(\E_\mu|Q-D_*|\le\sqrt a=\varepsilon/2\).
Now truncate after \(m=\lceil2H/\varepsilon\rceil\) queries, outputting
the conditional expectation of \(Q\) at each cutoff leaf. Both outputs
lie in \([0,1]\) and agree unless \(K>m\), so Markov's inequality gives
\[
 \E_\mu|Q-D|
 \le\E_\mu|Q-D_*|+\Pr(K>m)
 \le\varepsilon/2+H/m\le\varepsilon.
\]
This proves \eqref{ian:bridge-depth}.

A \(T\)-query quantum acceptance function is globally in \([0,1]\)
and has degree at most \(2T\), as used in Girish Theorem~4.13, p.~21;
take \(d=\max\{1,2T\}\). With \(z=1/\varepsilon\) and \(b=2T+1\), so
that \(1\le d/a\le4bz^2\), the depth \eqref{ian:bridge-depth} is at most
\[
 P(T,z)=1+2z\left[
 \frac{2b^2}{c_L}(4bz^2)^{\lceil C_L\rceil}
 +\frac{1+B_G(4bz^2)^{\lceil K_G\rceil}}{c_G}
       (4bz^2)^{\lceil C_G\rceil}
 \right],
\]
whose constants depend only on the six generalized-AA constants.
Since \(\Lambda\) is pinning-closed, the same bound holds from every
positive-probability pinning of every member, proving the hereditary
assertion.
\end{proof}

\begin{remark}[Relation to the existing reveal criterion]
\label{ian:bridge-versus-reveal}
Proposition~\ref{tht:reveal-prop} in
\path{foundations/transfers/02_reveal_progress.tex}
assumes polynomial reveal progress; the bridge also accepts influence
progress, which need not reveal anything immediately: under the
uniform two-spin law, \(Q=(1+x_1x_2)/2\) has \(I_1(Q)=I_2(Q)=1/4\),
whereas either singleton has zero reveal gain. Local queries are
charged to total influence and Global steps to conditional variance.
This two-potential amortisation adapts the recursion of AA14
Theorem~3.3, as presented in Girish Theorem~4.13, p.~21, and is not a
direct application of Proposition~\ref{tht:reveal-prop}.
\end{remark}
